% Persian print preamble. Loaded via include-in-header from _quarto.yml.
%
% ENGINE: LuaLaTeX

% Persian digits for every LaTeX counter: chapter, section, page, footnote.
\babelprovide[import, main, maparabic]{persian}

% Chinese font support via Babel character-based font switching
\babelprovide[import, onchar=ids fonts]{chinese}

\directlua{
  local function font_exists(name)
    local ok = pcall(function() return fonts.names.resolve(name) end)
    return ok
  end
  if font_exists("Songti SC") then
    tex.print([[\babelfont[chinese]{rm}{Songti SC}]])
  elseif font_exists("Noto Serif CJK SC") then
    tex.print([[\babelfont[chinese]{rm}{Noto Serif CJK SC}]])
  elseif font_exists("Noto Sans CJK SC") then
    tex.print([[\babelfont[chinese]{rm}{Noto Sans CJK SC}]])
  elseif font_exists("PingFang SC") then
    tex.print([[\babelfont[chinese]{rm}{PingFang SC}]])
  else
    tex.print([[\babelfont[chinese]{rm}{Hiragino Sans GB}]])
  end
}

% Pandoc loads selnolig unconditionally under LuaTeX. Turn it off.
\providecommand{\selnoligoff}{}
\selnoligoff

% Footnote rule on the right. LaTeX's \footnoterule is a raw \hrule pinned to
% the left of the column and babel does not mirror it. This box inherits
% \bodydir TRT, so its first item lands at the right edge: rule, then \hfill.
\renewcommand{\footnoterule}{%
  \kern-3pt\hbox to \columnwidth{\vrule width .4\columnwidth height .4pt\hfill}\kern2.6pt}

% Line spacing for Persian typography
\usepackage{setspace}
\setstretch{1.25}
